\documentclass[9pt]{article}
\usepackage[a4paper,margin=0.35in]{geometry}
\usepackage{amsmath,amssymb,amsfonts,mathtools}
\setlength{\jot}{4pt}
\usepackage{xcolor}
\setlength{\parindent}{0pt}
\setlength{\parskip}{5pt}
\def\mf#1{\mathbb{#1}}
\newcommand{\Leps}{L_{\epsilon}}
\newcommand{\KK}{\mathcal K(x-y)\!\cdot\!\mathcal K(x'-y')}
\newcommand{\Wyy}{\left[U^{ce'}(y') - U^{ce'}(x')\right] \left[U^{ce}(y) - U^{ce}(x)\right] \rho^{e'}(x')\, \rho^{e}(x)}
\newcommand{\Wzz}{\left[U^{ce'}(z) - U^{ce'}(x')\right]\left[U^{ce}(z) - U^{ce}(x)\right]\rho^{e'}(x')\,\rho^{e}(x)}
\title{Group II, Row IV: complete derivation, all steps\\
{\normalsize (both poles subtracted at the start; the endpoint term absorbed into the standard
$\mathcal P_{gg}$, so the finite terms are genuinely finite)}}
\author{}
\date{}
\begin{document}
\maketitle
\vspace{-2.0em}

{\color{red}\textbf{Why this version.} In the earlier grouping the ``finite terms'' contained
the endpoint entry $\propto\log(1-\frac{k^{+}}{\vee})\,\Leps$, which carries
$\frac{2}{\epsilon}$ --- i.e.\ it was \emph{not} finite. That entry is exactly the difference
between the locally subtracted plus prescription and the standard one
(Step~8 below). Here the plus prescription is converted to the standard $[0,1]$ convention, so
the endpoint term is absorbed into the textbook $\mathcal P_{gg}$ inside the DGLAP category,
where its pole belongs with the collinear kernel. After that, the finite block consists of the
boundary and azimuthal terms only, and both are finite in every sense: no $\frac1\epsilon$, no
$\mu$, no $\log\Lambda$, no large $\log\frac{\vee}{k^{+}}$.}

\textbf{Notation.} $\mathcal K^{m}(u)=\frac{u^{m}}{u^{2}}$,
$K(a,b,z)=\mathcal K(a-z)\!\cdot\!\mathcal K(b-z)$, $r=y'-y$,
$\Leps=\frac{2}{\epsilon}+2\gamma_{E}+\log\frac{r^{2}\mu^{2}_{\overline{MS}}}{4}$,
$\bar\xi=\frac{k^{+}}{\vee}$, $\xi_{\max}=\frac{k^{+}}{k^{+}+\Lambda}$,
$W(u,u')=[U^{ce'}(u')-U^{ce'}(x')][U^{ce}(u)-U^{ce}(x)]\rho^{e'}(x')\rho^{e}(x)$,
$\mathcal N=\frac{1}{(2\pi)^{3}}\frac{g^{4}N_{c}}{8\pi^{5}}\frac{1}{k^{+}}$,
$C=\pi\mathcal N=\frac{1}{(2\pi)^{3}}\frac{g^{4}N_{c}}{8\pi^{4}}\frac{1}{k^{+}}$.

\section{Step 1: the master expression}
\begin{align} \label{eq:first}
     \frac{d^{3}N_{\rm NLO}^{\rm}}{d^{3}k}  & = \frac{1}{(2\pi)^{3}}\,\frac{g^{4}}{4\pi^{4}}\,\int_{\Lambda}^{\vee-k^{+}}\, \frac{dp^{+}}{2\pi}\int_{z,w,w^{\prime},x'} e^{-ik(w^{\prime}-w)}\,  \int_{x,y,y^{\prime}}\,\frac{(w-z)^{m}\,(x-y)^{k}\,(w'-z)^{m'}\,(x'-y')^{k'}}{(w-z)^{2}\,(x-y)^{2}\,(w'-z)^{2}\,(x'-y')^{2}}\,\frac{p^{+}\,k^{+}}{(p^{+}+k^{+})^{4}}\notag \\ & \delta^{(2)}\left[y-w+\frac{p^{+}}{(p^{+}+k^{+})}\,(w-z)\right]\,\times\, \delta^{(2)}\left[y'-w'+\frac{p^{+}}{(p^{+}+k^{+})}\,(w'-z)\right]\times\notag \\ &  \left[\delta_{km}\,\delta_{ij}-\frac{(p^{+}+k^{+})}{p^{+}}\,\delta_{jm}\,\delta_{ik}-\frac{(p^{+}+k^{+})}{k^{+}}\,\delta_{im}\,\delta_{kj}\right]\times\, \left[\delta_{k'm'}\,\delta_{ij}-\frac{(p^{+}+k^{+})}{p^{+}}\,\delta_{jm'}\,\delta_{ik'}-\frac{(p^{+}+k^{+})}{k^{+}}\,\delta_{im'}\,\delta_{k'j}\right]\notag \\ & \left[f^{c'd'a}\,U^{c'e'}(y')\,U^{d'b}(z)\,\rho^{e'}(x')-f^{c'd'a}\,
    U^{d'b}(z)\,U^{c'e'}(x')\,\rho^{e'}(x')\right]\times \left[f^{cda}\,U^{ce}(y)\,U^{db}(z)\,\rho^{e}(x)-f^{cda}\,
    U^{db}(z)\,U^{ce}(x)\,\rho^{e}(x)\right]
\end{align}

\section{Step 2: colour algebra}
The matrix elements are numbers, so they can be reordered. Factoring the common
$U^{d'b}(z)U^{db}(z)$ out of the product of the two colour brackets,
\begin{align}
[\,\cdots'][\,\cdots]=f^{c'd'a}f^{cda}\;U^{d'b}(z)\,U^{db}(z)\;
\left[U^{c'e'}(y')-U^{c'e'}(x')\right]\left[U^{ce}(y)-U^{ce}(x)\right]\rho^{e'}(x')\rho^{e}(x)\,.
\end{align}
The adjoint Wilson lines are real orthogonal matrices, so summing over $b$:
$U^{d'b}(z)U^{db}(z)=(UU^{T})^{d'd}=\delta^{d'd}$. Then
$f^{c'd'a}f^{cd'a}=N_{c}\,\delta^{c'c}$, and the colour factor collapses to
\begin{align}
N_{c}\,\left[U^{ce'}(y')-U^{ce'}(x')\right]\left[U^{ce}(y)-U^{ce}(x)\right]\rho^{e'}(x')\rho^{e}(x)
=N_{c}\,W(y,y')\,.
\end{align}
All $z$-dependence of the colour structure is gone: $z$ now appears only in the kernels.

\section{Step 3: integrating the $\delta$-functions over $w,w'$}
Write $\lambda=\frac{p^{+}}{p^{+}+k^{+}}$, so $1-\lambda=\frac{k^{+}}{p^{+}+k^{+}}$. The
argument of the first $\delta$ is $y-(1-\lambda)w-\lambda z$, which vanishes at
$w=\frac{y-\lambda z}{1-\lambda}$, and $\frac{\partial(\text{arg})}{\partial w}=-(1-\lambda)\mathbb 1$ in
two dimensions, so
\begin{align}
\delta^{(2)}\big[y-w+\lambda(w-z)\big]
=\Big(\frac{p^{+}+k^{+}}{k^{+}}\Big)^{2}\,
\delta^{(2)}\Big[w-\frac{y-\lambda z}{1-\lambda}\Big]\,,
\end{align}
and the same for the primed leg. At the support of the $\delta$'s:
\begin{align}
w-z=\frac{y-z}{1-\lambda}=\frac{p^{+}+k^{+}}{k^{+}}\,(y-z)\,,\qquad
\frac{(w-z)^{m}}{(w-z)^{2}}=\frac{k^{+}}{p^{+}+k^{+}}\,\frac{(y-z)^{m}}{(y-z)^{2}}\,,\qquad
w'-w=\frac{p^{+}+k^{+}}{k^{+}}\,(y'-y)\,.
\end{align}
The prefactor assembles as
\begin{align}
\underbrace{\frac{p^{+}k^{+}}{(p^{+}+k^{+})^{4}}}_{\rm original}\times
\underbrace{\Big(\frac{p^{+}+k^{+}}{k^{+}}\Big)^{4}}_{\rm two\ Jacobians}\times
\underbrace{\Big(\frac{k^{+}}{p^{+}+k^{+}}\Big)^{2}}_{\rm two\ kernels}
=\frac{p^{+}}{k^{+}\,(p^{+}+k^{+})^{2}}\,,
\end{align}
and the phase becomes $e^{-ik\cdot(y'-y)\frac{p^{+}+k^{+}}{k^{+}}}$. So after Steps 2--3:
\begin{align}\label{eq:afterdelta}
\frac{d^{3}N_{\rm NLO}}{d^{3}k}
=\frac{1}{(2\pi)^{3}}\frac{g^{4}N_{c}}{4\pi^{4}}
\int_{\Lambda}^{\vee-k^{+}}\frac{dp^{+}}{2\pi}\,\frac{p^{+}}{k^{+}(p^{+}+k^{+})^{2}}
\int_{x,x',y,y',z} e^{-ik\cdot(y'-y)\frac{p^{+}+k^{+}}{k^{+}}}\,
\mathcal K^{m}(y-z)\mathcal K^{k}(x-y)\mathcal K^{m'}(y'-z)\mathcal K^{k'}(x'-y')\;
B^{mkm'k'}\;W(y,y')
\end{align}
with $B^{mkm'k'}$ the contraction of the two square brackets over $i,j$, done next.

\section{Step 4: the tensor contraction}
Write $A=\frac{p^{+}+k^{+}}{p^{+}}$, $B=\frac{p^{+}+k^{+}}{k^{+}}$. Multiplying out the two
brackets and summing over $i,j$ ($\delta_{ii}=d_{\perp}$), the nine terms are
\begin{center}
\begin{tabular}{llll}
$\delta\!\cdot\!\delta$: $d_{\perp}\delta_{km}\delta_{k'm'}$ &
$\delta\!\cdot\!(-A)$: $-A\,\delta_{km}\delta_{k'm'}$ &
$\delta\!\cdot\!(-B)$: $-B\,\delta_{km}\delta_{k'm'}$ & \\
$(-A)\!\cdot\!\delta$: $-A\,\delta_{km}\delta_{k'm'}$ &
$A^{2}$: $A^{2}\delta_{mm'}\delta_{kk'}$ &
$AB$: $AB\,\delta_{mk'}\delta_{km'}$ & \\
$(-B)\!\cdot\!\delta$: $-B\,\delta_{km}\delta_{k'm'}$ &
$BA$: $AB\,\delta_{mk'}\delta_{km'}$ &
$B^{2}$: $B^{2}\delta_{mm'}\delta_{kk'}$ &
\end{tabular}
\end{center}
\begin{align}
B^{mkm'k'}=\big[d_{\perp}-2A-2B\big]\,\delta_{km}\delta_{k'm'}
+\big[A^{2}+B^{2}\big]\,\delta_{mm'}\delta_{kk'}
+2AB\,\delta_{mk'}\delta_{km'}
\end{align}
(checked by explicit index summation). Multiplying by the weight
$\frac{p^{+}}{k^{+}(p^{+}+k^{+})^{2}}$ of \eqref{eq:afterdelta}:
\begin{align}
-2(A+B)\,\frac{p^{+}}{k^{+}(p^{+}+k^{+})^{2}}=-\frac{2}{k^{+2}}\,,\qquad
2AB\,\frac{p^{+}}{k^{+}(p^{+}+k^{+})^{2}}=+\frac{2}{k^{+2}}\,,\qquad
(A^{2}+B^{2})\,\frac{p^{+}}{k^{+}(p^{+}+k^{+})^{2}}
=\frac{1}{k^{+}}\Big(\frac{p^{+}}{k^{+2}}+\frac{1}{p^{+}}\Big)\,,
\end{align}
so the full weight is
\begin{align}\label{eq:weights}
\frac{1}{k^{+}}\Bigg[\frac{d_{\perp}\,p^{+}}{(p^{+}+k^{+})^{2}}\,\delta_{km}\delta_{k'm'}
+\frac{2}{k^{+}}\big(\delta_{mk'}\delta_{km'}-\delta_{km}\delta_{k'm'}\big)
+\Big(\frac{p^{+}}{k^{+2}}+\frac{1}{p^{+}}\Big)\,\delta_{mm'}\delta_{kk'}\Bigg]\,.
\end{align}

\textbf{The antisymmetric channel drops.} For every term whose colour factor is
$z$-independent (which is all of \eqref{eq:afterdelta}: the colour factor is $W(y,y')$), the
$z$-integral $\int_{z}\mathcal K^{m}(y-z)\mathcal K^{m'}(y'-z)$ is a symmetric tensor in
$(m,m')$, while $\delta_{mk'}\delta_{km'}-\delta_{km}\delta_{k'm'}$ contracted with
$\mathcal K^{k}(x-y)\mathcal K^{k'}(x'-y')$ is antisymmetric under $m\leftrightarrow m'$;
their contraction vanishes identically, at every $p^{+}$, before any expansion. It is dropped
from here on.

\section{Step 5: the $\xi$ substitution}
With $\xi=\frac{k^{+}}{p^{+}+k^{+}}$: $p^{+}=\frac{k^{+}(1-\xi)}{\xi}$,
$p^{+}+k^{+}=\frac{k^{+}}{\xi}$, $dp^{+}=-\frac{k^{+}}{\xi^{2}}d\xi$, and the limits
$p^{+}\in[\Lambda,\vee-k^{+}]$ map to $\xi:\xi_{\max}\to\bar\xi$ (descending); flipping to
ascending absorbs the minus sign of $dp^{+}$. The channel weights become
\begin{align}
\frac{p^{+}}{(p^{+}+k^{+})^{2}}=\frac{\xi(1-\xi)}{k^{+}}\,,\qquad
\frac{p^{+}}{k^{+2}}=\frac{1-\xi}{\xi\,k^{+}}\,,\qquad
\frac{1}{p^{+}}=\frac{\xi}{(1-\xi)\,k^{+}}\,,\qquad
e^{-ikr\frac{p^{+}+k^{+}}{k^{+}}}=e^{-ikr\frac{1}{\xi}}\,,
\end{align}
and the prefactor assembles to $\mathcal N$:
$\frac{g^{4}N_{c}}{4\pi^{4}}\cdot\frac{1}{2\pi}\cdot\frac{k^{+}}{\xi^{2}}\cdot\frac{1}{k^{+}}\cdot\frac{1}{k^{+}}
=\frac{g^{4}N_{c}}{8\pi^{5}}\frac{1}{k^{+}}\frac{1}{\xi^{2}}$. Contracting the kernels with
the two surviving tensor structures
($\delta_{km}\delta_{k'm'}$: $[\mathcal K(y-z)\!\cdot\!\mathcal K(x-y)][\mathcal K(y'-z)\!\cdot\!\mathcal K(x'-y')]$;
$\delta_{mm'}\delta_{kk'}$: $K(y,y',z)\,\KK$):
\begin{align}\label{eq:xiform}
\boxed{\;\frac{d^{3}N_{\rm NLO}}{d^{3}k}
=\mathcal N
\int_{\bar\xi}^{\xi_{\max}}\frac{d\xi}{\xi^{2}}
\int_{x,x',y,y',z} e^{-ik\cdot r\frac{1}{\xi}}
\Big\{d_{\perp}\,\xi(1-\xi)\big[\mathcal K(y-z)\!\cdot\!\mathcal K(x-y)\big]\big[\mathcal K(y'-z)\!\cdot\!\mathcal K(x'-y')\big]
+\Big[\frac{1-\xi}{\xi}+\frac{\xi}{1-\xi}\Big]K(y,y',z)\,\KK\Big\}W(y,y')\;}
\end{align}
This is the whole Row-IV cross-section, exactly, with the $z$-integral not yet performed.

\section{Step 6: subtract both poles, once each}
The curly bracket of \eqref{eq:xiform} is singular only at the two ends:
$\xi\to1$ ($p^{+}\to0$, soft target-side gluon, regulated by $\Lambda$) through
$\frac{\xi}{1-\xi}$, and $\xi\to0$ ($p^{+}\to\vee$, fast projectile-side gluon) through
$\frac{1-\xi}{\xi}$ combined with the measure and the phase (the position-space collapse
$y,y'\to z$, since the $\delta$-functions of \eqref{eq:first} give $y-z=\xi(w-z)$,
$y'-z=\xi(w'-z)$). Define one counterterm per pole, each the exact endpoint limit of its
channel:
\begin{align}
T_{J}(\xi)&=\frac{1}{1-\xi}\,e^{-ik\cdot r}\,K(y,y',z)\,\KK\;W(y,y')
&&(\text{the }\tfrac{\xi}{1-\xi}\text{ channel at }\xi=1\text{: }e^{-ikr/\xi}\to e^{-ikr},\ \tfrac1\xi\to1)\,,\label{eq:TJ}\\
T_{B}(\xi)&=\frac{1}{\xi^{2}}\,\frac{1-\xi}{\xi}\,e^{-ik\cdot r\frac{1}{\xi}}\,K(y,y',z)\,K(x,x',z)\;W(z,z)
&&(\text{the }\tfrac{1-\xi}{\xi}\text{ channel with }\KK\,W(y,y')\to K(x,x',z)W(z,z))\,.\label{eq:TB}
\end{align}
Then, exactly,
\begin{align}\label{eq:split}
\frac{d^{3}N_{\rm NLO}}{d^{3}k}
=\underbrace{\mathcal N\!\int_{\bar\xi}^{\xi_{\max}}\!\! d\xi\!\int\Big[\{\text{integrand of \eqref{eq:xiform}}\}-T_{J}-T_{B}\Big]}_{\textstyle d^{3}N_{\rm reg}}
+\underbrace{\mathcal N\!\int_{\bar\xi}^{\xi_{\max}}\!\! d\xi\!\int T_{J}}_{\textstyle J}
+\underbrace{\mathcal N\!\int_{\bar\xi}^{\xi_{\max}}\!\! d\xi\!\int T_{B}}_{\textstyle B}\;.
\end{align}
Each counterterm is subtracted once and added back once: nothing can be double counted.
Finiteness of $d^{3}N_{\rm reg}$ at the two ends: at $\xi\to1$ the soft channel minus $T_{J}$
carries $\frac{\frac1\xi e^{-ikr/\xi}-e^{-ikr}}{1-\xi}$, finite at $\xi=1$, so
$\xi_{\max}\to1$ at a cost $O(\Lambda)$; at $\xi\to0$ the $\frac{1-\xi}{\xi}$ channel minus
$T_{B}$ is
\begin{align}\label{eq:combo}
\frac{1}{\xi^{2}}\frac{1-\xi}{\xi}\,e^{-ikr\frac1\xi}K(y,y',z)
\Big[\KK\,W(y,y')-K(x,x',z)\,W(z,z)\Big],
\end{align}
and in the variables $y=z+\xi(w-z)$, $y'=z+\xi(w'-z)$ the square bracket vanishes linearly in
$\xi$ (because $y,y'\to z$ and $\mathcal K(x-z)\!\cdot\!\mathcal K(x'-z)=K(x,x',z)$
identically), so the $\bar\xi$ end produces no large logarithm. In $T_{B}$'s own integral the
upper limit may be taken to $1$ at a cost $O(\Lambda)$ (its integrand is regular at $\xi=1$).

\section{Step 7: the $\xi\to1$ counterterm $J$ --- JIMWLK}
$T_{J}$ depends on $\xi$ only through $\frac{1}{1-\xi}$, so the $\xi$-integral factorizes and
is elementary (antiderivative $-\log(1-\xi)$, substituted upper minus lower):
\begin{align}
\int_{\bar\xi}^{\xi_{\max}}\frac{d\xi}{1-\xi}
=\Big[-\log(1-\xi)\Big]_{\bar\xi}^{\xi_{\max}}
=\log\frac{1-\bar\xi}{1-\xi_{\max}}
=\log\frac{(\vee-k^{+})\,k^{+}}{\vee\,\Lambda}+O(\Lambda)
={\color{blue}\log\frac{k^{+}}{\Lambda}}+\log\Big(1-\frac{k^{+}}{\vee}\Big)+O(\Lambda)\,.
\end{align}
Hence
\begin{align}\label{eq:J}
J=\mathcal N\,\Big[{\color{blue}\log\frac{k^{+}}{\Lambda}}+\log\Big(1-\frac{k^{+}}{\vee}\Big)\Big]
\int_{x,x',y,y',z} e^{-ik\cdot r}\,K(y,y',z)\,\KK\,W(y,y')\,.
\end{align}
The $\log\frac{k^{+}}{\Lambda}$ part is the \textbf{JIMWLK term}, in its $z$-form:
\begin{align}\label{eq:jimwlk}
\boxed{\;\frac{d^{3}N_{\rm JIMWLK}}{d^{3}k}
=+\,\frac{1}{(2\pi)^{3}}\frac{g^{4}N_{c}}{8\pi^{5}}\frac{1}{k^{+}}\,{\color{blue}\log\frac{k^{+}}{\Lambda}}
\int_{x,y,z,x',y'}e^{-ik\cdot(y'-y)}\,K(y,y',z)\,\KK\,\Wyy\;}
\end{align}
(with $\frac1\pi\int_{z}K(y,y',z)=-\Leps$ this is the $\Leps$-form with the all-negative
bracket). The $\log(1-\frac{k^{+}}{\vee})$ part of \eqref{eq:J} --- call it the
\emph{endpoint piece} --- is \textbf{not} put into the finite block: Step~8 shows it is
exactly what converts the plus prescription to the standard one, and it is absorbed there.

\section{Step 8: the endpoint piece $=$ the plus-prescription conversion}
Inside $d^{3}N_{\rm reg}$, the soft channel appears in the \emph{locally subtracted} form.
Writing $h(\xi)\equiv\frac{1}{\xi}\,e^{-ikr\frac1\xi}\,K(y,y',z)\,\KK\,W(y,y')$ (so that the
$\frac{\xi}{1-\xi}$ channel over the measure is $\frac{h(\xi)}{1-\xi}$ and $T_{J}=\frac{h(1)}{1-\xi}$),
the standard plus distribution on $[0,1]$ obeys the exact identity
\begin{align}\label{eq:stdplus}
\int_{\bar\xi}^{1}d\xi\,\frac{h(\xi)}{(1-\xi)_{+}}
=\int_{\bar\xi}^{1}d\xi\,\frac{h(\xi)-h(1)}{1-\xi}
\;+\;h(1)\,\log\Big(1-\bar\xi\Big)\,.
\end{align}
The first term on the right is precisely the locally subtracted soft channel of
$d^{3}N_{\rm reg}$; the second is precisely the endpoint piece of \eqref{eq:J}
($\log(1-\bar\xi)=\log(1-\frac{k^{+}}{\vee})$, and $h(1)=e^{-ikr}K(y,y',z)\KK W$). Therefore
\begin{align}
(\text{locally subtracted soft channel})+(\text{endpoint piece})
=\mathcal N\int_{\bar\xi}^{1}\frac{d\xi}{\xi^{2}}\,\frac{\xi}{(1-\xi)_{+}}
\int e^{-ikr\frac1\xi}K(y,y',z)\,\KK\,W(y,y')
\end{align}
with the \emph{standard} plus prescription --- the one for which
$\mathcal P_{gg,+}(\xi)=2N_{c}\big[\xi(1-\xi)+\frac{1-\xi}{\xi}+\frac{\xi}{(1-\xi)_{+}}\big]$
is the textbook splitting function. From here on the soft channel means this standard-plus
form, and no endpoint term exists anywhere.

\section{Step 9: the $\xi\to0$ counterterm $B$ --- BFKL $+$ boundary}
$T_{B}$ depends on $\xi$ only through the phase and explicit powers. Change variables at
fixed $z$: $y=z+\xi(w-z)$, $y'=z+\xi(w'-z)$ (exact, invertible). Then, factor by factor:
(i) $d^{2}y\,d^{2}y'=\xi^{4}\,d^{2}w\,d^{2}w'$;
(ii) $y-z=\xi(w-z)$, $y'-z=\xi(w'-z)$, $y'-y=\xi(w'-w)$;
(iii) $K(y,y',z)=\frac{\xi^{2}(w-z)\cdot(w'-z)}{\xi^{2}(w-z)^{2}\,\xi^{2}(w'-z)^{2}}=\frac{K(w,w',z)}{\xi^{2}}$;
(iv) $e^{-ikr/\xi}=e^{-ik\xi(w'-w)/\xi}=e^{-ik(w'-w)}$ exactly;
(v) $K(x,x',z)$, $W(z,z)$ untouched. Hence the transverse integral of $T_{B}$ equals
$\xi^{4}\cdot\xi^{-2}\cdot\frac{1}{\xi^{2}}\cdot\frac{1-\xi}{\xi}$ wait --- keeping the
explicit powers of \eqref{eq:TB} outside: $\int_{y,y'}T_{B}
=\frac{1}{\xi^{2}}\frac{1-\xi}{\xi}\;\xi^{4}\,\frac{1}{\xi^{2}}\int_{w,w'}e^{-ik(w'-w)}K(w,w',z)K(x,x',z)W(z,z)
=\frac{1-\xi}{\xi}\times(\xi\text{-independent})$, so
\begin{align}
B=\mathcal N\Big[\int_{\bar\xi}^{1}d\xi\,\frac{1-\xi}{\xi}\Big]\mathcal I\,,\qquad
\int_{\bar\xi}^{1}d\xi\,\frac{1-\xi}{\xi}
=\Big[\log\xi-\xi\Big]_{\bar\xi}^{1}
={\color{blue}\log\frac{\vee}{k^{+}}}-\Big(1-\frac{k^{+}}{\vee}\Big)\,,\qquad
\mathcal I=\int_{x,x',z,w,w'}\!\!e^{-ik(w'-w)}K(w,w',z)K(x,x',z)\,\Wzz\,.
\end{align}
The $\log\frac{\vee}{k^{+}}$ part is the \textbf{BFKL term},
\begin{align}\label{eq:bfkl}
\boxed{\;\frac{d^{3}N_{\rm BFKL}}{d^{3}k}
=+\,\frac{1}{(2\pi)^{3}}\,\frac{g^{4}N_{c}}{8\pi^{5}}\,\frac{1}{k^{+}}\,{\color{blue}\log\frac{\vee}{k^{+}}}\;\mathcal I\;}
\end{align}
(UV finite and $\mu$-independent: every integral above was done in $d=2$ with no regulator),
and the remainder is the \textbf{boundary term}
\begin{align}\label{eq:boundary}
-\,\frac{1}{(2\pi)^{3}}\,\frac{g^{4}N_{c}}{8\pi^{5}}\,\frac{1}{k^{+}}\,\Big(1-\frac{k^{+}}{\vee}\Big)\,\mathcal I\,.
\end{align}

\section{Step 10: the doubly subtracted remainder --- DGLAP, RC, azimuthal}
After Steps 6 and 8, $d^{3}N_{\rm reg}$ (with the endpoint piece absorbed) contains three
structures: (a) the $d_{\perp}\xi(1-\xi)$ channel, (c) the standard-plus soft channel
$\frac{\xi}{(1-\xi)_{+}}$, and (p) the $\mu$-free pair \eqref{eq:combo}. The pair (p) is
finite in $d=2$ and is kept as it stands. In (a) and (c) the colour factor $W(y,y')$ is
$z$-independent, so the $z$-integral is performed in dimensional regularization
($d_{\perp}=2-\epsilon$, $\overline{MS}$):
\begin{align}\label{eq:zint}
\int_{z}\mathcal K^{m}(y-z)\,\mathcal K^{m'}(y'-z)
=-\frac{\pi}{2}\,\Leps\,\delta^{mm'}-\pi\,\frac{r^{m}r^{m'}}{r^{2}}\,.
\end{align}
\textbf{(a)} contract \eqref{eq:zint} with $\mathcal K^{m}(x-y)\mathcal K^{m'}(x'-y')$:
\begin{align}
d_{\perp}\,\xi(1-\xi)\Big[-\frac{\pi}{2}\Leps\,\KK-\pi\,\frac{(\mathcal K(x-y)\cdot r)(\mathcal K(x'-y')\cdot r)}{r^{2}}\Big].
\end{align}
Using $\frac{d_{\perp}}{2}\Leps=\big(1-\frac{\epsilon}{2}\big)\Leps=\Leps-1+O(\epsilon)$
(the $-\frac{\epsilon}{2}\cdot\frac{2}{\epsilon}=-1$), the $\KK$ piece is
$-\pi\,\xi(1-\xi)\,(\Leps-1)\,\KK$.
\textbf{(c)} here the contraction is the trace: $\delta^{mm'}$ against \eqref{eq:zint} gives
$\int_{z}K(y,y',z)=-\frac{\pi}{2}\Leps\,d_{\perp}-\pi=-\pi\Leps+O(\epsilon)$ (the
$+\frac{\epsilon\pi}{2}\Leps=+\pi$ cancels the $-\pi$ from the $\frac{rr}{r^{2}}$ trace), so
(c) gives $-\pi\,\frac{\xi}{(1-\xi)_{+}}\,\Leps\,\KK$.

Multiplying by $\mathcal N\frac{d\xi}{\xi^{2}}$ ($\mathcal N\pi=C$) and using
$\frac{\xi(1-\xi)}{\xi^{2}}=\frac{1-\xi}{\xi}$:
\begin{itemize}
\item the $\Leps$ pieces of (a) and (c):
$-C\int\frac{d\xi}{\xi^{2}}\big[\xi(1-\xi)+\frac{\xi}{(1-\xi)_{+}}\big]\Leps\,\KK
=-\frac{g^{2}}{8\pi^{2}}\int\frac{d\xi}{\xi^{2}}\big[\mathcal P_{gg,+}(\xi)-2N_{c}\tfrac{1-\xi}{\xi}\big]\Leps\,(\dots)$,
since $2N_{c}[\xi(1-\xi)+\frac{\xi}{(1-\xi)_{+}}]=\mathcal P_{gg,+}-2N_{c}\frac{1-\xi}{\xi}$
and $C=\frac{g^{2}}{8\pi^{2}}\,2N_{c}\,\frac{1}{(2\pi)^{3}}\frac{g^{2}}{2\pi^{2}}\frac{1}{k^{+}}$;
\item the ``$-1$'' leftover of (a), $+C\frac{1-\xi}{\xi}\KK$, and the r-tensor of (a),
$-C\,d_{\perp}\frac{1-\xi}{\xi}\frac{(\mathcal K\cdot r)(\mathcal K'\cdot r)}{r^{2}}$,
combine into the \textbf{azimuthal term} ($d_{\perp}\to2$, no pole multiplies it).
\end{itemize}
Finally add and subtract $\frac{\beta_{0}}{2}\delta(1-\xi)$ inside the bracket
($\mathcal P_{gg}=\mathcal P_{gg,+}+\frac{\beta_{0}}{2}\delta(1-\xi)$,
$\beta_{0}=\frac{11N_{c}}{3}$): the $\delta$ sits at $\xi=1$, where $\frac{1}{\xi^{2}}=1$ and
the phase is the LO phase, producing the \textbf{RC term}. Altogether:
\begin{align}\label{eq:reg}
d^{3}N_{\rm reg}+(\text{endpoint piece})
=&\underbrace{-\,\frac{1}{(2\pi)^{3}}\,\frac{g^{2}}{2\pi^{2}}\,\frac{g^{2}}{8\pi^{2}}\,\frac{1}{k^{+}}
\int_{\bar\xi}^{1}\frac{d\xi}{\xi^{2}}\Big[\mathcal P_{gg}(\xi)-2N_{c}\frac{1-\xi}{\xi}\Big]
\int_{x,y,x',y'}e^{-ik\cdot r\frac{1}{\xi}}\,\KK\;\Leps\;\Wyy}_{\textstyle\rm DGLAP\ (non\text{-}soft),\ standard\ plus\ prescription}\notag\\
&\underbrace{+\,\frac{1}{(2\pi)^{3}}\frac{g^{4}N_{c}}{8\pi^{5}}\frac{1}{k^{+}}
\int_{\bar\xi}^{1}\frac{d\xi}{\xi^{2}}\,\frac{1-\xi}{\xi}
\int_{x,y,x',y',z}e^{-ik\cdot r\frac{1}{\xi}}K(y,y',z)
\Big[\KK\,W(y,y')-K(x,x',z)\,W(z,z)\Big]}_{\textstyle\rm DGLAP\ (soft\ pair):\ free\ of\ \frac{2}{\epsilon},\ \mu\ and\ \log\frac{\vee}{k^{+}}}\notag\\
&\underbrace{+\,\frac{1}{(2\pi)^{3}}\,\frac{g^{2}}{2\pi^{2}}\,\frac{g^{2}\beta_{0}}{16\pi^{2}}\,\frac{1}{k^{+}}
\int_{x,y,x',y'}e^{-ik\cdot r}\,\KK\;\Leps\;\Wyy}_{\textstyle\rm RC}\notag\\
&\underbrace{+\,\frac{1}{(2\pi)^{3}}\,\frac{g^{4}N_{c}}{8\pi^{4}}\,\frac{1}{k^{+}}
\int_{\bar\xi}^{1}d\xi\,\frac{1-\xi}{\xi}
\int_{x,y,x',y'}e^{-ik\cdot r\frac{1}{\xi}}
\Big[\KK-2\,\frac{(\mathcal K(x-y)\!\cdot\!r)(\mathcal K(x'-y')\!\cdot\!r)}{r^{2}}\Big]\Wyy}_{\textstyle\rm azimuthal}\,.
\end{align}

\section{Step 11: result, and why the finite terms are finite}
\begin{align}
\frac{d^{3}N_{\rm NLO}}{d^{3}k}
=\underbrace{\eqref{eq:jimwlk}}_{\rm JIMWLK}
+\underbrace{\eqref{eq:bfkl}}_{\rm BFKL}
+\underbrace{\rm DGLAP\ \eqref{eq:reg}_{1,2}}_{\rm collinear}
+\underbrace{\rm RC\ \eqref{eq:reg}_{3}}_{\rm coupling}
+\underbrace{\rm boundary\ \eqref{eq:boundary}\;+\;azimuthal\ \eqref{eq:reg}_{4}}_{\textstyle\rm finite\ terms}\,.
\end{align}
\textbf{The finite block now consists of the boundary and azimuthal terms only, and both are
finite in every sense:}
\begin{itemize}
\item[(i)] \emph{No ultraviolet logarithms.} Neither contains $\frac{1}{\epsilon}$,
$\gamma_{E}$ or $\mu$: the boundary term's transverse integral $\mathcal I$ converges in
$d=2$ (at $z\to x,x'$ the kernels are $\sim\frac{1}{|z-x|}$, integrable; at large $z$ the
Wilson-line differences cut it off), and the azimuthal term multiplies no pole, so
$d_{\perp}=2$ there exactly.
\item[(ii)] \emph{No $\log\Lambda$.} All $\Lambda$-dependence left with
\eqref{eq:jimwlk}, by Step 7.
\item[(iii)] \emph{No large $\log\frac{\vee}{k^{+}}$.} The boundary coefficient
$(1-\frac{k^{+}}{\vee})$ is bounded by $1$. The azimuthal $\xi$-integral converges as
$\bar\xi\to0$: at fixed positions, $u=\frac1\xi$ turns it into
$\int^{\vee/k^{+}}du(\frac1u-\frac{1}{u^{2}})e^{-ik\cdot ru}$, oscillatory-convergent; and the
collapse region $r=O(\xi)$ is doubly suppressed, by the measure ($\xi^{2}$) and by the exact
vanishing of the angular average
$\int_{0}^{2\pi}\frac{d\theta_{\hat r}}{2\pi}[\mathcal K\!\cdot\!\mathcal K'-2(\mathcal K\!\cdot\!\hat r)(\mathcal K'\!\cdot\!\hat r)]=0$.
\end{itemize}
The earlier ``endpoint term'' $\propto\log(1-\frac{k^{+}}{\vee})\Leps$ no longer appears: by
Step 8 it is the exact difference between the locally subtracted and standard plus
prescriptions, and with the textbook $\mathcal P_{gg}$ its $\frac{2}{\epsilon}$ sits where it
belongs, inside the collinear DGLAP pole. The pole/log inventory of the row is therefore:
$\log\frac{k^{+}}{\Lambda}$ only in JIMWLK; $\log\frac{\vee}{k^{+}}$ only in BFKL (both with
UV-finite, $\mu$-free coefficients); $\frac{2}{\epsilon}$ only in the non-soft DGLAP kernel
and the RC term (collinear physics and coupling renormalization); and a finite block that is
genuinely finite.
\end{document}
